%%%%%%%%%%%%%%%%%%%%%%%%
%
% $Authors: Gautam Ramesh, Siddanth Sharma, Shubhankar Dumka $
% $Date: 2026-01-27 $
% $Path: /Manual/Assembly/Hints.tex $
% $Version: 3.1 (Clean Final Version) $
%
%%%%%%%%%%%%%%%%%%%%%%%%

\chapter{General Assembly Hints}

Before beginning the physical assembly of the License Plate Detection System, please review the following critical guidelines. These instructions ensure a safe, efficient, and error-free setup process for the Jetson Nano hardware and its peripheral ecosystem.

\section{Preparation and Safety Measures}

\begin{itemize}
	\item \textbf{Work Environment:} Assemble all components on a clean, dry, and non-conductive surface with proper ventilation.
	\item \textbf{Electrostatic Discharge (ESD):} Handle the reComputer J1010 and camera only by their edges. Use an antistatic wrist strap or ground yourself before touching any internal electronics.
	\item \textbf{Cable Integrity:} Arrange cables neatly without tension. Avoid crossing power and data cables where possible to minimize interference.
	\item \textbf{Power Protocol:} Use only the official \SI{5}{\volt}/\SI{4}{\ampere} adapter. Always disconnect the wall plug before adding or removing internal components.
\end{itemize}

% --- ADVANCED TIKZ SYSTEM READINESS HUD ---
\begin{figure}[H]
	\centering
	\begin{tikzpicture}[
		font=\sffamily\small,
		node distance=0.5cm,
		% Styles
		hud_box/.style={rectangle, rounded corners=5pt, draw=blue!60, fill=blue!5, minimum width=4cm, minimum height=1.5cm, align=left, thick},
		header/.style={font=\sffamily\bfseries\color{blue!80!black}},
		icon/.style={circle, fill=blue!60, text=white, minimum size=0.6cm, font=\bfseries\footnotesize}
		]
		% Nodes
		\node[hud_box] (elec) {\textbf{Electrical Safety}\\ \footnotesize 5V/4A Standard\\ \footnotesize No Hot-Plugging};
		\node[hud_box, right=of elec] (env) {\textbf{Environment}\\ \footnotesize ESD-Safe Zone\\ \footnotesize Passive Airflow};
		\node[hud_box, right=of env] (conn) {\textbf{Connectivity}\\ \footnotesize USB 3.0 Mapping\\ \footnotesize No Cable Tension};
		
		% Connectors
		\draw[thick, blue!30] (elec.east) -- (env.west);
		\draw[thick, blue!30] (env.east) -- (conn.west);
		
		% Visual Indicators
		\node[icon, above=0.1cm of elec] {1};
		\node[icon, above=0.1cm of env] {2};
		\node[icon, above=0.1cm of conn] {3};
	\end{tikzpicture}
	\caption{Pre-assembly Readiness Dashboard: Ensuring hardware safety standards are met before initialization.}
\end{figure}

\section{Assembly Preparation Checklist}
Adherence to the following checklist is mandatory for first-time system verification.

\begin{table}[H]
	\centering
	\caption{Assembly Preparation Checklist (Ramesh et al., 2026)}
	\begin{tabular}{|L{4cm}|L{9cm}|}
		\hline
		\textbf{Item} & \textbf{Description / Action Required} \\ \hline
		Work Surface & Clean, dry, and non-conductive area \\ \hline
		Power Adapter & Official \SI{5}{\volt} - \SI{4}{\ampere} Type-C Supply \\ \hline
		SD Card & \SI{32}{GB} (or larger) Class 10 media \\ \hline
		Camera Module & Microsoft LifeCam Show (USB 3.0 compatible) \\ \hline
		Static Safety & ESD wrist strap or grounding point nearby \\ \hline
	\end{tabular}
\end{table}

\section{Post-Assembly Verification}
After completing physical connections, verify the following system states:
\begin{enumerate}
	\item The Green Power LED on the reComputer chassis illuminates steadily.
	\item The HDMI display output is visible on the monitor.
	\item The USB camera is detected and functioning properly.
	\item No abnormal heat or noise occurs during operation.
\end{enumerate}

\vspace{0.4cm}
\noindent
Following these general assembly hints ensures reliable installation and provides a stable foundation for further software configuration.